\documentclass[12pt, a4paper]{article}

% ── Packages ──────────────────────────────────────────────────────────────────
\usepackage[utf8]{inputenc}
\usepackage[T1]{fontenc}
\usepackage[english]{babel}
\usepackage{geometry}
\geometry{a4paper, margin=2.5cm}

\usepackage{amsmath, amssymb}
\usepackage{booktabs}
\usepackage{array}
\usepackage{longtable}
\usepackage{multirow}
\usepackage{xcolor}
\usepackage{graphicx}
\usepackage{hyperref}
\usepackage{parskip}
\usepackage{titlesec}
\usepackage{fancyhdr}
\usepackage{enumitem}
\usepackage{caption}
\usepackage{lmodern}
\usepackage{microtype}
\usepackage{tcolorbox}
\tcbuselibrary{skins}

% ── Colours ───────────────────────────────────────────────────────────────────
\definecolor{ugblue}{HTML}{1a3a5c}
\definecolor{lightblue}{HTML}{2e6da4}
\definecolor{rowgray}{HTML}{f5f5f5}

% ── Hyperref ──────────────────────────────────────────────────────────────────
\hypersetup{
    colorlinks=true,
    linkcolor=lightblue,
    urlcolor=lightblue,
    citecolor=lightblue,
    pdftitle={Reorganisation of an Outpatient Radiology Department},
    pdfauthor={Simulation Modelling and Analysis 2025-2026},
}

% ── Section formatting ────────────────────────────────────────────────────────
\titleformat{\section}{\large\bfseries\color{ugblue}}{\thesection.}{0.6em}{}[\vspace{-0.3em}\textcolor{ugblue}{\rule{\linewidth}{0.8pt}}]
\titleformat{\subsection}{\normalsize\bfseries\color{lightblue}}{\thesubsection.}{0.5em}{}
\titleformat{\subsubsection}{\normalsize\itshape\color{lightblue}}{\thesubsubsection.}{0.5em}{}

% ── Header / footer ───────────────────────────────────────────────────────────
\pagestyle{fancy}
\fancyhf{}
\fancyhead[L]{\small\color{gray}Reorganisation of an Outpatient Radiology Department}
\fancyhead[R]{\small\color{gray}SMA 2025--2026}
\fancyfoot[C]{\thepage}
\renewcommand{\headrulewidth}{0.4pt}

% ── Recommendation box ────────────────────────────────────────────────────────
\newtcolorbox{recbox}{
    enhanced,
    colback=ugblue!8,
    colframe=ugblue,
    boxrule=1pt,
    arc=3pt,
    left=8pt, right=8pt, top=6pt, bottom=6pt,
    fontupper=\itshape\color{ugblue},
}

% ── Table row colours ─────────────────────────────────────────────────────────
\usepackage{colortbl}
\newcommand{\headerrow}{\rowcolor{ugblue}\color{white}}
\newcommand{\altrow}{\rowcolor{rowgray}}

% ─────────────────────────────────────────────────────────────────────────────
\begin{document}

% ── Cover page ────────────────────────────────────────────────────────────────
\begin{titlepage}
    \centering
    \vspace*{2cm}
    {\fontsize{28}{34}\selectfont\bfseries\color{ugblue}
        Reorganisation of an\\[4pt] Outpatient Radiology\\[4pt] Department\par}
    \vspace{0.8cm}
    \textcolor{ugblue}{\rule{0.7\linewidth}{1.5pt}}
    \vspace{0.6cm}

    {\large Simulation Modelling and Analysis (F000941)}\\[4pt]
    {\large Academic Year 2025--2026}\\[4pt]
    {\large Prof.\ Broos Maenhout}

    \vspace{1.2cm}
    {\large \textbf{Group [GROEPSNUMMER]}}\\[0.5cm]
    \begin{tabular}{ll}
        [Naam 1] & [Studentnummer] \\
        [Naam 2] & [Studentnummer] \\
        [Naam 3] & [Studentnummer] \\
        [Naam 4] & [Studentnummer] \\
        [Naam 5] & [Studentnummer] \\
        [Naam 6] & [Studentnummer] \\
    \end{tabular}

    \vfill
    {\normalsize Ghent University --- Master of Science in Management of Operations}
\end{titlepage}

\tableofcontents
\newpage

% ─────────────────────────────────────────────────────────────────────────────
\section{Introduction}
% ─────────────────────────────────────────────────────────────────────────────

\subsection{System Description}

This study considers the redesign of the cyclic appointment schedule for a single-server outpatient radiology department. The department serves two patient types: \textbf{elective} patients, who make appointments in advance, and \textbf{urgent} patients, who arrive without an appointment and must be treated on the same day.

\textbf{Schedule structure.} The planning horizon repeats weekly. The department is open Monday through Saturday; Thursday and Saturday are half-day sessions. Each day is divided into time slots of 15~minutes — the expected scan duration for an elective patient. The current schedule contains 146 elective slots and 14 urgent slots per week. Resources (nurses, doctors) are assumed to be always available.

\textbf{Elective patients.} Elective patients call the front desk (open Monday--Friday, 08:00--17:00) to request an appointment. Inter-arrival times between phone calls follow a memoryless negative-exponential distribution, with the number of daily calls following a Poisson distribution ($\lambda_e = 28.345$ patients/day). Patients are assigned to the next available elective slot on a first-come, first-served basis. At the moment of the appointment, patients may arrive early or late: unpunctuality follows a $\mathcal{N}(0,\, 2.5~\text{min})$ distribution. The probability of a no-show is 2\%. Elective scan durations follow a $\mathcal{N}(15~\text{min},\, 3~\text{min})$ distribution.

\textbf{Urgent patients.} Urgent patients arrive unannounced during opening hours according to a negative-exponential inter-arrival distribution, with $\lambda_u = 2.5$ patients per full day and $\lambda_u = 1.25$ per half day. They are assigned to the next available urgent slot. If no slot remains on that day, they are served in overtime after all elective patients have been treated. Urgent patients have variable scan durations depending on scan type; Table~\ref{tab:urgent_scan} summarises the distribution.

\begin{table}[ht]
\centering
\caption{Scan duration distribution for urgent patients by scan type.}
\label{tab:urgent_scan}
\small
\begin{tabular}{lccc}
\toprule
\headerrow \textbf{Scan type} & \textbf{Frequency} & \textbf{Mean (min)} & \textbf{Std.\ dev.\ (min)} \\
\midrule
Brain              & 70\% & 15.0 & 2.5 \\
\altrow Spine --- Lumbar vertebrae  & 10\% & 17.5 & 1.0 \\
Spine --- Cervical vertebrae & 10\% & 22.5 & 2.5 \\
\altrow Abdomen MRI         &  5\% & 30.0 & 1.0 \\
Others             &  5\% & 30.0 & 4.5 \\
\bottomrule
\end{tabular}
\end{table}

\textbf{Performance measures.} Four performance measures are tracked: (1)~elective appointment waiting time ($Wt_e^{app}$): time between phone call and appointment; (2)~elective scan waiting time ($Wt_e^{scan}$): time between physical arrival and scan start; (3)~urgent scan waiting time ($Wt_u^{scan}$): time between urgent patient arrival and scan start; (4)~daily overtime. The primary objective is to minimise the weighted sum:
\[
    \text{Minimise} \quad w_e \cdot \overline{Wt_e^{app}} + w_u \cdot \overline{Wt_u^{scan}}
    \quad \text{with} \quad w_e = \tfrac{1}{168}, \quad w_u = \tfrac{1}{9}
\]
Elective scan waiting time and overtime serve as \textit{secondary} objectives, used only to break ties between statistically equivalent configurations.

\subsection{Research Hypotheses}

To evaluate the impact of the three design factors on system performance, we formulate the following research hypotheses:

\textbf{Main hypothesis:}
\begin{itemize}
    \item $H_{0_1}$: The current appointment schedule (Strategy~1, 14 urgent slots, Rule~1) minimises the primary objective function; no alternative configuration performs significantly better.
    \item $H_{A_1}$: At least one alternative configuration exists that achieves a significantly lower weighted waiting time than the current schedule.
\end{itemize}

\textbf{Effect of timing strategy:}
\begin{itemize}
    \item $H_{0_2}$: The timing strategy does not significantly affect the primary objective value; all three strategies produce equivalent weighted waiting times.
    \item $H_{A_2}$: At least one timing strategy achieves a significantly lower primary objective value. Specifically, distributing urgent slots more evenly throughout the session (Strategies~2 and~3) is expected to reduce both urgent scan waiting time and elective appointment waiting time compared to end-of-block placement (Strategy~1).
\end{itemize}

\textbf{Effect of urgent slot capacity:}
\begin{itemize}
    \item $H_{0_3}$: The number of urgent slots does not significantly affect the primary objective; performance is equivalent across the range of 10--17 slots.
    \item $H_{A_3}$: There exists an optimal number of urgent slots that minimises the primary objective. Configurations with too few slots increase urgent scan waiting times; configurations with too many slots displace elective appointments and increase elective waiting times.
    \item $H_{0_4}$: Configurations with 18 or more urgent slots are operationally stable and comparable to configurations with fewer slots.
    \item $H_{A_4}$: Configurations with 18 or more urgent slots are systemically unstable due to capacity saturation and must be excluded from the experimental design.
\end{itemize}

\textbf{Effect of appointment scheduling rule:}
\begin{itemize}
    \item $H_{0_5}$: The appointment scheduling rule does not significantly affect the primary objective value.
    \item $H_{A_5}$: At least one appointment rule achieves a significantly different primary objective value. The rules are expected to primarily affect the secondary objective (elective scan waiting time) through their impact on punctuality-induced congestion, with negligible effect on the primary objective.
\end{itemize}

\subsection{Overview of Research Test Design}

To test the above hypotheses, we adopt a \textbf{full-factorial screening design} across three factors:

\begin{table}[ht]
\centering
\caption{Factors and levels in the experimental design.}
\label{tab:factors}
\small
\begin{tabular}{>{\bfseries}p{3.2cm} p{9.8cm} c}
\toprule
\headerrow \textbf{Factor} & \textbf{Levels} & \textbf{\#} \\
\midrule
Timing Strategy &
  \textit{Strategy 1}: urgent slots placed at end of each morning/afternoon block \newline
  \textit{Strategy 2}: urgent slots evenly distributed over the day \newline
  \textit{Strategy 3}: one urgent slot after every 6 consecutive elective slots & 3 \\
\altrow Urgent Slot Capacity & 10, 11, 12, 13, 14, 15, 16, 17 slots per week & 8 \\
Appointment Rule &
  \textit{Rule 1}: Plain FCFS \newline
  \textit{Rule 2}: Bailey-Welch ($K=2$) \newline
  \textit{Rule 3}: Blocking ($B=2$) \newline
  \textit{Rule 4}: Benchmarking ($k_\sigma = 0.5$) & 4 \\
\bottomrule
\end{tabular}
\end{table}

This yields $3 \times 8 \times 4 = \mathbf{96}$ configurations. A full-factorial design is chosen over a fractional design because the total number of configurations is computationally tractable ($\leq$100) and because we wish to detect potential interaction effects between strategy and slot count without aliasing.

Configurations with 18 or more urgent slots are excluded from the design: preliminary analysis using antithetic variates revealed that these configurations are systemically unstable (see Section~\ref{sec:variance}).

The analysis proceeds through the following phases: (1)~steady-state determination via Welch's method, (2)~independent sampling via the Batch Means Method, (3)~variance reduction via Antithetic Variables, (4)~full-factorial screening of all 96 designs, and (5)~comparative analysis and tiebreaking using secondary objectives.

% ─────────────────────────────────────────────────────────────────────────────
\section{Statistical Analysis \& Warm-up}
% ─────────────────────────────────────────────────────────────────────────────

\subsection{Welch's Method and Warm-up Period}

Before applying the Batch Means Method, it is essential to determine and remove the initial warm-up period to eliminate transient initialization bias. We applied Welch's method across 10 distinct representative experimental designs, chosen to ensure a diverse mix of urgent slots, scheduling rules, and timing strategies.

By plotting the moving average of the objective function for these designs, we visually identified the point at which output stabilized. Configurations operating at near-full capacity (such as those with 18 urgent slots) were excluded from this baseline determination, as their extreme instability skews the steady-state convergence.

Based on the worst-case convergence observed, we identified a cutoff of approximately 63~weeks. To maintain robustness across the majority of our experiments, we conservatively selected a universal warm-up period of $\mathbf{W_{\text{warmup}} = 60}$~\textbf{weeks}. This removes initialization bias for approximately 90\% of our tested scenarios.

\subsection{Batch Means Method}

To conduct a statistically valid steady-state analysis, we applied the Batch Means Method. Because consecutive simulation outputs are autocorrelated, the batch means method groups observations into batches to form approximately independent data points.

Based on worst-case autocorrelation analysis, we determined:
\begin{itemize}
    \item \textbf{Batch size} $M = 80$ weeks (lag-based correlation falls below 0.01)
    \item \textbf{Number of batches} $L = 22$ (sufficient for target precision)
\end{itemize}

This yields a total simulation run length of:
\[
    W_{\text{total}} = W_{\text{warmup}} + L \times M = 60 + 22 \times 80 = 1820 \text{ weeks}
\]

To establish a baseline of natural system variance, we first simulated 11 representative experimental designs. Table~\ref{tab:bmm_baseline} summarises the 95\% confidence intervals and precision metrics.

\begin{table}[ht]
\centering
\caption{Batch Means Method results for 11 representative designs (standard run).}
\label{tab:bmm_baseline}
\small
\begin{tabular}{lccccccc}
\toprule
\headerrow
\textbf{Design} & $\bar{D}$ & $S^2$ & $S$ & $\epsilon$ & \textbf{95\% CI} & \textbf{CV} & $\gamma$ \\
\midrule
S1-U14-R1 & 0.4783 & 0.0011 & 0.0339 & 0.0150 & [0.4633 ; 0.4933] & 0.0709 & 0.0314 \\
\altrow S1-U18-R4 & 0.9140 & 0.0659 & 0.2567 & 0.1138 & [0.8002 ; 1.0278] & 0.2809 & 0.1245 \\
S3-U17-R2 & 0.5838 & 0.0214 & 0.1462 & 0.0648 & [0.5190 ; 0.6486] & 0.2503 & 0.1110 \\
\altrow S1-U11-R3 & 0.4514 & 0.0003 & 0.0167 & 0.0074 & [0.4440 ; 0.4588] & 0.0369 & 0.0164 \\
S2-U17-R2 & 0.6527 & 0.0220 & 0.1484 & 0.0658 & [0.5869 ; 0.7185] & 0.2274 & 0.1008 \\
\altrow S3-U14-R3 & 0.4487 & 0.0011 & 0.0326 & 0.0145 & [0.4342 ; 0.4631] & 0.0727 & 0.0322 \\
S3-U17-R3 & 0.5863 & 0.0215 & 0.1466 & 0.0650 & [0.5214 ; 0.6513] & 0.2500 & 0.1108 \\
\altrow S3-U13-R3 & 0.4317 & 0.0006 & 0.0242 & 0.0107 & [0.4209 ; 0.4424] & 0.0562 & 0.0249 \\
S1-U14-R4 & 0.4769 & 0.0012 & 0.0341 & 0.0151 & [0.4618 ; 0.4920] & 0.0714 & 0.0317 \\
\altrow S2-U16-R1 & 0.5787 & 0.0059 & 0.0768 & 0.0341 & [0.5446 ; 0.6127] & 0.1328 & 0.0589 \\
S2-U13-R4 & 0.5233 & 0.0007 & 0.0262 & 0.0116 & [0.5117 ; 0.5349] & 0.0500 & 0.0222 \\
\bottomrule
\end{tabular}
\end{table}

The results indicate that while some configurations (e.g., S1-U11-R3, $\gamma = 1.64\%$) are highly stable, others (e.g., S1-U18-R4, $\gamma = 12.45\%$) exhibit substantial natural variance, highlighting the necessity of variance reduction techniques.

% ─────────────────────────────────────────────────────────────────────────────
\section{Variance Reduction}
\label{sec:variance}
% ─────────────────────────────────────────────────────────────────────────────

To obtain more accurate estimates without increasing computational burden, we implemented variance reduction using \textbf{Antithetic Variables}. For every standard simulation batch generated using random number sequence $U$, we generated an antithetic counterpart using $1-U$. These paired outputs are negatively correlated. By averaging both runs:
\[
    Y_k = \frac{X_k + X'_k}{2}
\]
the overall variance is substantially reduced. Table~\ref{tab:bmm_antithetic} presents the results for the same 11 representative designs after applying antithetic variates.

\begin{table}[ht]
\centering
\caption{Batch Means Method results for 11 representative designs (antithetic run, $Y_k = (X_k + X'_k)/2$).}
\label{tab:bmm_antithetic}
\small
\begin{tabular}{lccccccc}
\toprule
\headerrow
\textbf{Design} & $\bar{Y}$ & $S^2_Y$ & $S_Y$ & $\epsilon$ & \textbf{95\% CI} & \textbf{CV} & $\gamma$ \\
\midrule
S1-U14-R1  & 0.4880 & 0.0006 & 0.0240 & 0.0106 & [0.4774 ; 0.4986] & 0.0491 & 0.0218 \\
\altrow S1-U18-R4  & 2.3740 & 0.8888 & 0.9428 & 0.4180 & [1.9560 ; 2.7920] & 0.3971 & 0.1761 \\
S3-U17-R2  & 0.7391 & 0.0509 & 0.2256 & 0.1000 & [0.6391 ; 0.8391] & 0.3052 & 0.1353 \\
\altrow S1-U11-R3  & 0.4565 & 0.0001 & 0.0085 & 0.0038 & [0.4528 ; 0.4603] & 0.0186 & 0.0082 \\
S2-U17-R2  & 0.8084 & 0.0515 & 0.2268 & 0.1006 & [0.7078 ; 0.9090] & 0.2806 & 0.1244 \\
\altrow S3-U14-R3  & 0.4593 & 0.0006 & 0.0240 & 0.0106 & [0.4486 ; 0.4699] & 0.0523 & 0.0232 \\
S3-U17-R3  & 0.7418 & 0.0510 & 0.2259 & 0.1001 & [0.6417 ; 0.8420] & 0.3045 & 0.1350 \\
\altrow S3-U13-R3  & 0.4390 & 0.0002 & 0.0155 & 0.0069 & [0.4321 ; 0.4459] & 0.0353 & 0.0157 \\
S1-U14-R4  & 0.4865 & 0.0006 & 0.0241 & 0.0107 & [0.4759 ; 0.4972] & 0.0495 & 0.0219 \\
\altrow S2-U16-R1  & 0.6272 & 0.0075 & 0.0867 & 0.0384 & [0.5888 ; 0.6657] & 0.1382 & 0.0613 \\
S2-U13-R4  & 0.5298 & 0.0003 & 0.0164 & 0.0073 & [0.5225 ; 0.5371] & 0.0310 & 0.0138 \\
\bottomrule
\end{tabular}
\end{table}

\subsection{Impact on Stable Systems}

For configurations not pushing the system to capacity limits (11--14 urgent slots), antithetic pairing successfully halved the variance. For example:
\begin{itemize}
    \item \textbf{S1-U11-R3}: variance dropped from 0.0003 to 0.0001; $\gamma$ improved from 1.64\% to 0.82\%.
    \item \textbf{S1-U14-R1}: variance dropped from 0.0011 to 0.0006; $\gamma$ improved from 3.14\% to 2.18\%.
\end{itemize}

\subsection{Heavily Loaded Systems and the Exclusion of High Urgent Slot Counts}

The variance reduction technique exposed a critical vulnerability in heavily-loaded configurations. For \textbf{S1-U18-R4}, the standard mean was 0.9140. When the antithetic run was executed, the combined mean skyrocketed to 2.3740, with variance increasing rather than decreasing.

This occurs because queue waiting times are \textit{convex} relative to system utilization. When the antithetic sequence generates shorter inter-arrival times and longer service times simultaneously, an already saturated queue breaks down entirely, driving waiting times toward infinity. This is empirical proof that configurations with 18 or more urgent slots are \textbf{fundamentally unstable}. Consequently, urgent slot counts of 18, 19, and 20 were excluded from the full experimental design.

\subsection{Confidence Interval Estimation}

Since we estimate the population mean from $n = L = 22$ batches, we use the $t$-distribution:
\[
    CI = \left[\bar{Y} - t_{\alpha/2,\, n-1} \cdot \frac{S_Y}{\sqrt{n}} \;;\; \bar{Y} + t_{\alpha/2,\, n-1} \cdot \frac{S_Y}{\sqrt{n}}\right]
\]

For design \textbf{S1-U14-R1} after antithetic reduction ($\bar{Y} = 0.4880$, $S_Y = 0.0240$, $t_{0.025,21} \approx 2.0796$):
\[
    CI = \left[0.4880 - 2.0796 \cdot \frac{0.0240}{\sqrt{22}} \;;\; 0.4880 + 2.0796 \cdot \frac{0.0240}{\sqrt{22}}\right] = [0.4774 \;;\; 0.4986]
\]

% ─────────────────────────────────────────────────────────────────────────────
\section{Experimental Design (Screening)}
% ─────────────────────────────────────────────────────────────────────────────

\subsection{Objective Function}

The primary objective function to be minimized is:
\[
    \text{Minimize} \quad w_e \cdot \overline{Wt_e} + w_u \cdot \overline{Wt_u}
\]
where:
\begin{itemize}
    \item $w_e = 1/168$: weight for elective \textit{appointment} waiting time (normalized by 1 week = 168 hours)
    \item $w_u = 1/9$: weight for urgent \textit{scan} waiting time (normalized by 9 hours)
    \item $\overline{Wt_e}$: average elective appointment waiting time (hours)
    \item $\overline{Wt_u}$: average urgent scan waiting time (hours)
\end{itemize}

Elective scan waiting time and overtime are \textit{secondary} objectives, used only for tiebreaking.

\subsection{Full-Factorial Design}

We conducted a full-factorial screening design covering:
\begin{itemize}
    \item \textbf{3 timing strategies}: Strategy~1 (end-of-block), Strategy~2 (evenly distributed), Strategy~3 (after every 6 elective slots)
    \item \textbf{8 urgent slot counts}: 10, 11, 12, 13, 14, 15, 16, 17
    \item \textbf{4 appointment scheduling rules}: Rule~1 (Plain FCFS), Rule~2 (Bailey-Welch, $K=2$), Rule~3 (Blocking, $B=2$), Rule~4 (Benchmarking, $k_\sigma=0.5$)
\end{itemize}

This yields $3 \times 8 \times 4 = \mathbf{96}$ designs. Each was evaluated using the Batch Means Method with antithetic variates ($W_{\text{warmup}}=60$, $M=80$, $L=22$, seed~$=42$).

\subsection{Screening Results}

Table~\ref{tab:top20} presents the top 20 designs ranked by average objective value $\bar{Y}$. The complete results for all 96 designs are available in \texttt{full\_sweep\_antithetic.csv}.

\begin{table}[ht]
\centering
\caption{Top 20 designs from the full-factorial screening (antithetic BMM, $L=22$, $M=80$).}
\label{tab:top20}
\small
\begin{tabular}{clccccccc}
\toprule
\headerrow
\textbf{Rank} & \textbf{Design} & \textbf{Strat.} & \textbf{Urg.} & \textbf{Rule} & $\bar{Y}$ & $\epsilon$ & \textbf{95\% CI} \\
\midrule
1  & S3-U11-R4 & 3 & 11 & 4 & 0.4311 & 0.0040 & [0.4271 ; 0.4352] \\
\altrow 2  & S3-U11-R2 & 3 & 11 & 2 & 0.4321 & 0.0040 & [0.4280 ; 0.4361] \\
3  & S3-U11-R1 & 3 & 11 & 1 & 0.4325 & 0.0040 & [0.4284 ; 0.4365] \\
\altrow 4  & S3-U11-R3 & 3 & 11 & 3 & 0.4344 & 0.0041 & [0.4304 ; 0.4385] \\
5  & S3-U13-R4 & 3 & 13 & 4 & 0.4357 & 0.0069 & [0.4288 ; 0.4425] \\
\altrow 6  & S3-U13-R2 & 3 & 13 & 2 & 0.4366 & 0.0068 & [0.4297 ; 0.4434] \\
7  & S3-U13-R1 & 3 & 13 & 1 & 0.4371 & 0.0068 & [0.4302 ; 0.4439] \\
\altrow 8  & S3-U12-R4 & 3 & 12 & 4 & 0.4378 & 0.0050 & [0.4328 ; 0.4428] \\
9  & S3-U12-R2 & 3 & 12 & 2 & 0.4389 & 0.0050 & [0.4339 ; 0.4440] \\
\altrow 10 & S3-U13-R3 & 3 & 13 & 3 & 0.4390 & 0.0069 & [0.4321 ; 0.4459] \\
11 & S3-U12-R1 & 3 & 12 & 1 & 0.4392 & 0.0050 & [0.4342 ; 0.4441] \\
\altrow 12 & S3-U12-R3 & 3 & 12 & 3 & 0.4413 & 0.0051 & [0.4362 ; 0.4464] \\
13 & S3-U10-R4 & 3 & 10 & 4 & 0.4427 & 0.0035 & [0.4392 ; 0.4462] \\
\altrow 14 & S3-U10-R2 & 3 & 10 & 2 & 0.4438 & 0.0035 & [0.4402 ; 0.4473] \\
15 & S3-U10-R1 & 3 & 10 & 1 & 0.4440 & 0.0035 & [0.4405 ; 0.4476] \\
\altrow 16 & S3-U10-R3 & 3 & 10 & 3 & 0.4461 & 0.0036 & [0.4425 ; 0.4496] \\
17 & S1-U11-R2 & 1 & 11 & 2 & 0.4529 & 0.0037 & [0.4491 ; 0.4566] \\
\altrow 18 & S1-U11-R4 & 1 & 11 & 4 & 0.4532 & 0.0037 & [0.4495 ; 0.4570] \\
19 & S1-U11-R1 & 1 & 11 & 1 & 0.4547 & 0.0037 & [0.4510 ; 0.4585] \\
\altrow 20 & S3-U14-R4 & 3 & 14 & 4 & 0.4559 & 0.0106 & [0.4453 ; 0.4665] \\
\bottomrule
\end{tabular}
\end{table}

\subsection{Key Observations}

\subsubsection{Observation 1 — Strategy 3 dominates}

The entire top 16 positions are occupied by Strategy~3 designs. The first non-Strategy-3 design appears only at rank 17 (S1-U11-R2, $\bar{Y} = 0.4529$). Strategy~2 is clearly the worst-performing strategy.

\begin{table}[ht]
\centering
\caption{Best and worst objective value per strategy.}
\label{tab:per_strategy}
\small
\begin{tabular}{clcc}
\toprule
\headerrow \textbf{Strategy} & \textbf{Best Design} & \textbf{Best OV} & \textbf{Worst OV} \\
\midrule
1 & S1-U11-R2 & 0.4529 & 0.7883 \\
\altrow 2 & S2-U13-R4 & 0.5298 & 0.8113 \\
3 & S3-U11-R4 & 0.4311 & 0.7418 \\
\bottomrule
\end{tabular}
\end{table}

\subsubsection{Observation 2 — 11 urgent slots is the optimal capacity allocation}

Within Strategy~3, the best objective values are consistently achieved with 11 urgent slots. The OV increases as the number of urgent slots deviates from this optimum in both directions: too few slots lead to longer urgent scan waiting times, while too many slots crowd out elective appointments. Configurations with 15 or more urgent slots exhibit rapidly increasing variance, confirming the approach towards instability.

\begin{table}[ht]
\centering
\caption{Effect of urgent slot count on OV for Strategy 3.}
\label{tab:urgent_effect}
\small
\begin{tabular}{ccc}
\toprule
\headerrow \textbf{Urgent Slots} & \textbf{Best OV (S3)} & \textbf{Variance level} \\
\midrule
10 & 0.4427 & Very low \\
\altrow 11 & \textbf{0.4311} & Very low \\
12 & 0.4378 & Low \\
\altrow 13 & 0.4357 & Low \\
14 & 0.4559 & Medium \\
\altrow 15 & 0.4782 & Medium-high \\
16 & 0.5507 & High \\
\altrow 17 & 0.7383 & Very high \\
\bottomrule
\end{tabular}
\end{table}

\subsubsection{Observation 3 — Appointment rules have minimal impact}

Within any given strategy and urgent slot count, the four appointment scheduling rules produce very similar objective values. For the best configuration (S3-U11), the OV ranges only from 0.4311 (Rule~4) to 0.4344 (Rule~3) --- a spread of just 0.0033.

\begin{table}[ht]
\centering
\caption{Best overall design per appointment rule.}
\label{tab:per_rule}
\small
\begin{tabular}{clcc}
\toprule
\headerrow \textbf{Rule} & \textbf{Description} & \textbf{Best Design} & \textbf{Best OV} \\
\midrule
1 & Plain FCFS      & S3-U11-R1 & 0.4325 \\
\altrow 2 & Bailey-Welch ($K=2$) & S3-U11-R2 & 0.4321 \\
3 & Blocking ($B=2$) & S3-U11-R3 & 0.4344 \\
\altrow 4 & Benchmarking ($k_\sigma=0.5$) & S3-U11-R4 & 0.4311 \\
\bottomrule
\end{tabular}
\end{table}

% ─────────────────────────────────────────────────────────────────────────────
\section{Comparative Analysis}
% ─────────────────────────────────────────────────────────────────────────────

\subsection{Statistical Comparison of Top Designs}

To determine whether differences between top-performing designs are statistically significant, we examine whether their 95\% confidence intervals overlap. Two designs are considered significantly different if their confidence intervals do not overlap.

Comparing the top 4 designs (all S3-U11, Table~\ref{tab:top4_ci}):

\begin{table}[ht]
\centering
\caption{Confidence intervals for the top 4 designs.}
\label{tab:top4_ci}
\small
\begin{tabular}{lcccc}
\toprule
\headerrow \textbf{Design} & $\bar{Y}$ & $\epsilon$ & \textbf{CI lower} & \textbf{CI upper} \\
\midrule
S3-U11-R4 & 0.4311 & 0.0040 & 0.4271 & 0.4352 \\
\altrow S3-U11-R2 & 0.4321 & 0.4040 & 0.4280 & 0.4361 \\
S3-U11-R1 & 0.4325 & 0.0040 & 0.4284 & 0.4365 \\
\altrow S3-U11-R3 & 0.4344 & 0.0041 & 0.4304 & 0.4385 \\
\bottomrule
\end{tabular}
\end{table}

All four confidence intervals overlap substantially. We therefore \textbf{cannot statistically distinguish} between these four designs based on the current sample. All S3-U11 configurations perform equivalently well on the primary objective.

\subsection{Comparison Between Strategies}

The performance gap between Strategy~3 and the alternatives is \textbf{statistically significant}:

\begin{itemize}
    \item Comparing the best S3 design (S3-U11-R4, CI: $[0.4271;\, 0.4352]$) against the best S1 design (S1-U11-R2, CI: $[0.4491;\, 0.4566]$): the entire CI of S3 lies below the entire CI of S1, confirming that Strategy~3 is statistically superior to Strategy~1.
    \item The superiority over Strategy~2 is even clearer: the best S2 design (S2-U13-R4, CI: $[0.5225;\, 0.5371]$) lies entirely above the entire top-16 cluster of Strategy~3 designs.
\end{itemize}

\subsection{Recommended Configuration}

\begin{recbox}
Recommended configuration: \textbf{Strategy 3 with 11 urgent slots} (any appointment rule). \\[4pt]
If a single rule must be selected: \textbf{Rule 4 (Benchmarking)} achieves the lowest point estimate ($\bar{Y} = 0.4311$).
\end{recbox}

This recommendation represents a change from the current schedule (Strategy~1, 14 urgent slots, Rule~1), which scores $\bar{Y} \approx 0.488$ --- approximately \textbf{11\% worse} than the best design.

\subsection{Secondary Objectives}

The assignment specifies that elective scan waiting time and overtime serve as secondary objectives for tiebreaking. Since the top-performing S3-U11 designs are statistically indistinguishable on the primary objective, a formal tiebreak using secondary objectives could be applied in future work.

% ─────────────────────────────────────────────────────────────────────────────
\section{Conclusion}
% ─────────────────────────────────────────────────────────────────────────────

This study evaluated 96 appointment schedule configurations for an outpatient radiology department, covering three timing strategies, eight urgent slot counts (10--17), and four appointment rules. A statistically rigorous simulation framework was applied, including Welch's warm-up analysis, the Batch Means Method ($L=22$, $M=80$), and variance reduction via antithetic variates.

The main conclusions are:

\begin{enumerate}
    \item \textbf{Strategy 3 is clearly the best timing strategy.} Placing an urgent slot after every 6 consecutive elective slots produces consistently lower waiting times than end-of-block placement (Strategy~1) or even distribution (Strategy~2). The difference is statistically significant.

    \item \textbf{11 urgent slots is the optimal capacity allocation.} This number minimizes the weighted combination of elective appointment waiting time and urgent scan waiting time. Fewer slots increase urgent patient delays; more slots create appointment backlog for elective patients. Configurations with 16 or more urgent slots exhibit high instability.

    \item \textbf{Appointment rules have negligible impact on the primary objective.} All four rules perform nearly identically within the optimal S3-U11 configuration. The confidence intervals of the top 4 designs overlap, confirming no statistically significant difference.

    \item \textbf{The recommended configuration is S3-U11-R4} (Benchmarking rule), with an estimated objective value of $\bar{Y} = 0.4311$ --- an \textbf{11\% improvement} over the current schedule.
\end{enumerate}

Next steps to complete the analysis include: formal tiebreaking on secondary objectives (elective scan waiting time and overtime) for the top designs, and optionally a Ranking \& Selection procedure (e.g., Dudewicz-Dalal) to formally confirm the best design with a probability-of-correct-selection guarantee.

\end{document}
